\chapter{Overview Mata Kuliah dan Alur Rekayasa}
\label{ch:overview}

\ifajar
\section*{Tujuan Pembelajaran}
\begin{itemize}
  \item Menjelaskan beda mata kuliah ini dengan pemrograman web pengantar,
        yaitu pergeseran dari aplikasi yang berjalan menjadi aplikasi yang
        diukur, beserta peta empat belas pertemuannya.
  \item Memeriksa kesiapan lingkungan kerja untuk keempat belas bab dengan
        satu \textit{smoke test}, lalu menjelaskan sebab tiap pemeriksaan yang
        gagal.
  \item Merancang alur kerja tim berupa monorepo, cabang berumur pendek,
        \textit{code review}, dan Continuous Integration (CI) yang menolak
        perubahan cacat, berdasarkan hasil pengukuran.
\end{itemize}
\fi

\section{Pendahuluan}

Mata kuliah pemrograman web pengantar bertanya apakah aplikasinya berjalan.
Mata kuliah ini bertanya seberapa baik aplikasi itu berjalan: seberapa cepat,
seberapa tahan ketika sebagian sistemnya gagal, seberapa aman, dan bagaimana
membuktikan semuanya. Jawabannya selalu berupa angka yang diukur sendiri,
bukan kesan.

Bab ini adalah orientasi. Isinya peta keempat belas pertemuan, perkakas yang
dipakai beserta alasannya, dan alur kerja rekayasa yang berlaku sepanjang
semester. Bab ini juga menyiapkan tiga hal yang dipakai bab-bab berikutnya:
lingkungan kerja yang sudah diperiksa, alur Git dengan cabang berumur pendek,
dan kerangka Application Programming Interface (API) kecil yang dijaga CI.

Referensi utamanya adalah Winters dkk.\ untuk alur rekayasa pada tim yang
besar \cite{winters2020swe}, Chacon dan Straub untuk Git
\cite{chacon2014progit}, Humble dan Farley untuk CI \cite{humble2010cd}, dan
Fowler untuk pola \textit{repository} \cite{fowler2002poeaa}.

\section{Cakupan dan Posisi Mata Kuliah}

Kata "lanjut" pada mata kuliah ini tidak berarti fitur yang lebih rumit. Yang
berubah adalah pertanyaannya. Fitur yang sama, misalnya katalog toko buku,
dinilai dari sifat yang tidak terlihat di layar: cepatnya ketika diakses
ratusan orang, perilakunya ketika basis datanya lambat, ketahanannya terhadap
serangan, dan kemudahannya diubah tanpa merusak bagian lain.

Analoginya seperti beda menyetir mobil sendiri dengan mengemudikan bus kampus.
Keduanya sama-sama sampai tujuan. Namun bus harus tepat waktu setiap hari,
mengangkut banyak penumpang sekaligus, dan tetap jalan ketika satu bannya
bocor, sehingga sopirnya dinilai dari catatan perjalanan, bukan dari satu
perjalanan yang lancar.

Tabel \ref{tab:posisi01} merangkum pergeseran itu.

\begin{table}[!htb]
  \centering
  \caption{Pertanyaan pemrograman web pengantar lawan mata kuliah ini}
  \label{tab:posisi01}
  \begin{tabularx}{\textwidth}{lXX}
    \hline
    \textbf{Aspek} & \textbf{Pengantar} & \textbf{Mata kuliah ini} \\
    \hline
    pertanyaan utama & apakah fiturnya berjalan & seberapa baik berjalan, dan buktinya \\
    \hline
    bukti & demo di depan kelas & angka hasil pengukuran, misalnya p95 \\
    \hline
    lingkungan & satu laptop, satu proses & \textit{container}, beberapa layanan, beban \\
    \hline
    kegagalan & dihindari & disengaja, diukur, lalu diredam \\
    \hline
    mutu kode & lulus bila berjalan & diuji otomatis dan dijaga CI \\
    \hline
  \end{tabularx}
\end{table}

Persentil 50 dan 95, ditulis p50 dan p95, muncul di hampir setiap bab. Nilai
p50 adalah waktu yang tidak dilampaui separuh permintaan, sedangkan p95 adalah
waktu yang tidak dilampaui 95 dari tiap 100 permintaan. Yang kedua lebih
jujur, karena pengguna yang paling lama menunggu justru yang paling sering
pergi.

Dua contoh dari bab-bab berikutnya memperlihatkan jenis jawaban yang dicari:
\begin{enumerate}
  \item Bab 13 mengirim sepuluh peristiwa yang sama dengan tiga cara.
        \textit{Polling} berjeda setengah detik menyerahkannya dalam 199,8
        milidetik, sedangkan sambungan yang mendorong menyerahkannya dalam
        1,8 milidetik. Keduanya "berjalan", tetapi angkanya berbeda seratus
        kali lipat.
  \item Bab 14 menambah replika aplikasi dari satu menjadi delapan.
        \textit{Throughput} naik dari 106,8 menjadi 198,0 permintaan per
        detik lalu berhenti, karena basis data di belakangnya penuh. Tanpa
        pengukuran, empat replika tambahan itu terlihat seperti perbaikan.
\end{enumerate}

Karena itu tiap bab ditutup dengan Latihan yang menghasilkan angka. Angka
tiap mahasiswa akan berbeda, sedangkan yang dinilai adalah pola
perbandingannya dan penjelasan atas kejanggalannya.

\section{Peta Empat Belas Pertemuan}

Keempat belas pertemuan dibagi menjadi empat bagian. Bagian pertama
menyiapkan alat ukur. Bagian kedua menyusun lapisan data, \textit{cache}, dan
performa. Bagian ketiga berjalan berselang-seling di antara tiga jalur:
pengujian dan operasi, keamanan, dan arsitektur \textit{frontend}. Bagian
keempat menutup dengan sistem asinkron dan penskalaan. Gambar
\ref{fig:peta01} memperlihatkan urutannya.

\begin{figure}[!htb]
  \centering
  \begin{tikzpicture}[
    font=\scriptsize,
    kotak/.style={draw, rounded corners, minimum height=1.5cm,
                  text width=2.25cm, align=center, font=\scriptsize\bfseries},
    panah/.style={-{Stealth[length=2mm]}, thick}
  ]
    \node[kotak, fill=blue!8] (bagian1) at (0,0) {Fondasi dan pengukuran\\{\tiny\mdseries pertemuan 1--2}};
    \node[kotak, fill=blue!8] (bagian2) at (3.1,0) {Data, \textit{cache}, dan performa\\{\tiny\mdseries pertemuan 3--5}};
    \node[kotak, fill=orange!12] (bagian3) at (6.2,0) {Tiga jalur berselang\\{\tiny\mdseries pertemuan 6--12}};
    \node[kotak, fill=green!10] (bagian4) at (9.3,0) {Asinkron dan skala\\{\tiny\mdseries pertemuan 13--14}};

    \draw[panah] (bagian1) -- (bagian2);
    \draw[panah] (bagian2) -- (bagian3);
    \draw[panah] (bagian3) -- (bagian4);
    \node[font=\tiny, gray] at (3.1,-1.1) {topik 1 siap};
    \node[font=\tiny, gray, align=center] at (6.2,-1.2) {pengujian, keamanan,\\\textit{frontend} bergantian};
  \end{tikzpicture}
  \caption{Empat bagian dari keempat belas pertemuan}
  \label{fig:peta01}
\end{figure}

Seluruh pertemuan memakai satu aplikasi yang sama, yaitu toko buku kampus.
Mulai Bab 3, basis datanya berisi 20 kategori, 50.000 buku, 20.000 pelanggan,
200.000 pesanan, dan 600.000 baris pesanan, sehingga \textit{query} yang
lambat benar-benar terasa lambat. Bab-bab sesudahnya menambahkan
\textit{cache}, pengujian, keamanan, identitas pengguna, kontrak API, fitur
\textit{real-time}, dan rencana penskalaannya.

Urutan pertemuan disusun untuk menopang enam topik penelitian. Materi yang
paling dibutuhkan diletakkan paling awal, sehingga sebagian topik sudah dapat
dimulai sebelum semester selesai. Tabel \ref{tab:topik01} memuat keenam topik
beserta pertemuan yang menopangnya.

\begin{table}[!htb]
  \centering
  \caption{Enam topik penelitian dan pertemuan yang menopangnya}
  \label{tab:topik01}
  \begin{tabularx}{\textwidth}{lXll}
    \hline
    \textbf{Topik} & \textbf{Masalah yang diteliti} & \textbf{Pertemuan} & \textbf{Siap sesudah} \\
    \hline
    T1 & invalidasi \textit{cache} hibrida & 3, 4, 5 & pertemuan 5 \\
    \hline
    T2 & pencabutan token di \textit{gateway} & 7, 10, 12 & pertemuan 12 \\
    \hline
    T3 & transport \textit{real-time} adaptif & 5, 13 & pertemuan 13 \\
    \hline
    T4 & \textit{semantic cache} untuk Large Language Model (LLM) & 4, 5, 14 & pertemuan 14 \\
    \hline
    T5 & hidrasi selektif & 5, 8, 11 & pertemuan 11 \\
    \hline
    T6 & \textit{gate} kinerja di CI & 3, 6, 9 & pertemuan 9 \\
    \hline
  \end{tabularx}
\end{table}

Keenam topik memakai pola argumen yang sama: satu masalah, satu solusi yang
menimbulkan efek samping, lalu satu komponen yang ditambahkan untuk menekan
efek samping itu. Rincian tiap topik ada di \texttt{README.md} akar
repositori.

\section{Tech Stack dan Alasan Pemilihannya}

Perkakas dipilih dengan satu prinsip: yang terukur harus mekanismenya, bukan
\textit{framework}-nya. Karena itu sebagian besar server dan alat ukur ditulis
dengan pustaka standar Python, termasuk server HyperText Transfer Protocol
(HTTP)-nya. Server WebSocket di Bab 13, misalnya, ditulis
tanpa pustaka tambahan, sehingga jabat tangan dan bentuk \textit{frame}-nya
terlihat langsung di kode. Tabel \ref{tab:stack01} memuat perkakas lainnya.

\begin{table}[!htb]
  \centering
  \caption{Perkakas yang dipakai beserta alasan pemilihannya}
  \label{tab:stack01}
  \begin{tabularx}{\textwidth}{llX}
    \hline
    \textbf{Kebutuhan} & \textbf{Pilihan} & \textbf{Alasan} \\
    \hline
    server dan alat ukur & Python 3 & kode terbaca, pustaka standarnya cukup untuk HTTP \\
    \hline
    kode di browser & JavaScript, TypeScript & memang berjalan di browser, tipenya diperiksa di Bab 12 \\
    \hline
    basis data & PostgreSQL 16 & rencana \textit{query} dan \textit{isolation level} terlihat jelas \\
    \hline
    \textit{cache} & Redis 7 & \textit{cache} di luar proses yang dipakai bersama \\
    \hline
    lingkungan & Docker Compose & basis data yang sama di tiap laptop \\
    \hline
    pengukuran browser & Google Chrome, Playwright & membuka halaman tanpa jendela, lalu diukur \\
    \hline
    versi & Git, satu monorepo & naskah, slide, dan kode berubah bersama \\
    \hline
  \end{tabularx}
\end{table}

Perkakas sebanyak itu harus diperiksa sebelum praktikum pertama, bukan ketika
praktikum Bab 5 sudah berjalan setengah. Pemeriksaan singkat seperti itu
disebut \textit{smoke test}: belum menguji apa pun secara mendalam, hanya
memastikan tidak ada yang langsung berasap begitu dinyalakan.

Analoginya seperti daftar periksa di laboratorium kimia. Sebelum percobaan
dimulai, asisten memeriksa gas, air, dan alat ukur satu per satu. Alat yang
rusak ketahuan dalam lima menit, bukan di tengah percobaan yang bahannya sudah
telanjur dicampur.

\textit{Smoke test} bab ini memeriksa versi tiap perkakas, apakah Docker
benar-benar dapat dihubungi, apakah pustaka basis data sudah ada di venv
bersama, dan apakah port yang dipakai bab-bab berikutnya masih kosong.
Pemeriksaan port tertulis pada Listing \ref{lst:port01}.

\begin{lstlisting}[style=PythonStyle, caption={Memeriksa satu port tanpa mengirim data, tersedia di \berkas{sources/chapter01/periksa-lingkungan.py}}, label={lst:port01}]
def port_terpakai(port):
  """Mencoba memesan satu port di seluruh alamat, lalu langsung melepasnya.

  SO_REUSEADDR dipasang agar sambungan lama yang sedang ditutup tidak
  dianggap pemakai. Yang gagal dipesan hanya port yang sedang didengarkan
  proses lain.
  """
  with socket.socket(socket.AF_INET, socket.SOCK_STREAM) as soket:
    soket.setsockopt(socket.SOL_SOCKET, socket.SO_REUSEADDR, 1)
    try:
      soket.bind(("", port))
    except OSError:
      return True
  return False
\end{lstlisting}

Tabel \ref{tab:lingkungan01} memuat hasilnya pada laptop yang dipakai menulis
modul ini.

\begin{table}[!htb]
  \centering
  \caption{Hasil \textit{smoke test} lingkungan, hasil pengukuran nyata}
  \label{tab:lingkungan01}
  \begin{tabularx}{\textwidth}{XlX}
    \hline
    \textbf{Pemeriksaan} & \textbf{Hasil} & \textbf{Keterangan} \\
    \hline
    Python 3.11 atau lebih baru & lolos & 3.15.2 \\
    \hline
    Node.js 20 atau lebih baru & lolos & v22.23.4 \\
    \hline
    npm 10 atau lebih baru & lolos & 10.9.9 \\
    \hline
    Git 2.30 atau lebih baru & lolos & 2.43.0 \\
    \hline
    Docker 24 atau lebih baru & lolos & server 29.9.3 \\
    \hline
    Docker Compose 2 atau lebih baru & lolos & 5.7.0 \\
    \hline
    Chrome atau Chromium & lolos & 157.0.8103.38 \\
    \hline
    psycopg di \texttt{sources/.venv} & lolos & 3.3.5 \\
    \hline
    port bab 1 sampai 14 & gagal & 3 dari 37 terpakai: 5433, 8000, 8080 \\
    \hline
  \end{tabularx}
\end{table}

Baris terakhir memuat kejanggalan. Seluruh perkakas lolos, tetapi tiga port
sudah dipakai proyek lain di laptop yang sama: basis data di port 5433, server
aplikasi di 8000, dan server file statis di 8080. Ketiganya port milik Bab 3,
Bab 4, dan Bab 11. Tanpa pemeriksaan ini, kegagalannya baru muncul berminggu-minggu
kemudian, sebagai pesan \textit{address already in use} di tengah praktikum
bab yang bersangkutan.

Hasil itu punya dua konsekuensi praktis. Pertama, versi minimum ditulis
eksplisit di skripnya, sehingga "sudah terpasang" berarti sama bagi seluruh
mahasiswa. Kedua, satu hasil merah cukup untuk menahan praktikum, sama
seperti CI yang menahan penggabungan kode: masalahnya diselesaikan lebih
dahulu, bukan diakali di tengah jalan.

\section{Alur Kerja Rekayasa: Monorepo, Branching, dan Code Review}

Monorepo berarti seluruh bagian sistem disimpan dalam satu repositori Git.
Repositori mata kuliah ini sendiri berbentuk begitu, seperti pada Gambar
\ref{fig:monorepo01}: naskah, slide, dan kode tiap bab tinggal bersama,
sehingga satu perubahan pada kode dan pada naskah yang mengutipnya masuk
dalam satu commit. Winters dkk.\ mencatat manfaat yang sama pada skala jauh
lebih besar: satu versi untuk tiap pustaka, dan perubahan lintas proyek yang
dapat dilakukan sekaligus \cite{winters2020swe}.

\begin{figure}[!htb]
  \centering
  \begin{tikzpicture}[
    font=\scriptsize,
    kotak/.style={draw, rounded corners, minimum height=0.75cm,
                  minimum width=2.3cm, align=center, font=\scriptsize\bfseries},
    panah/.style={-{Stealth[length=2mm]}, thick}
  ]
    \node[kotak, fill=praditagreen!15, minimum width=4.2cm] (akar) at (4.2,1.4) {advanced-web-engineering};
    \node[kotak, fill=blue!8] (module) at (0,0) {module/\\{\tiny\mdseries naskah 14 bab}};
    \node[kotak, fill=blue!8] (slides) at (2.8,0) {slides/\\{\tiny\mdseries 14 sesi}};
    \node[kotak, fill=orange!12] (sources) at (5.6,0) {sources/\\{\tiny\mdseries chapter01--14}};
    \node[kotak, fill=blue!8] (figures) at (8.4,0) {figures/\\{\tiny\mdseries gambar bersama}};
    \foreach \anak in {module, slides, sources, figures} {
      \draw[thick] (akar.south) -- (\anak.north);
    }
    \node[font=\tiny, gray] at (5.6,-0.65) {kode tiap bab, satu CI};
  \end{tikzpicture}
  \caption{Repositori mata kuliah ini sebagai satu monorepo}
  \label{fig:monorepo01}
\end{figure}

Analoginya seperti satu lemari arsip fakultas dengan laci untuk tiap program
studi, alih-alih tiap dosen menyimpan arsipnya di rumah masing-masing. Ketika
aturan akademik berubah, seluruh laci diperbarui sekaligus di satu tempat,
dan tidak ada \textit{copy} aturan lama yang tertinggal di rumah seseorang.

Monorepo tidak menentukan cara bekerja di dalamnya. Pertanyaan berikutnya:
berapa lama sebuah cabang boleh hidup sebelum digabung ke \texttt{main}? Pada
\textit{trunk-based development}, cabang hidup paling lama beberapa hari, dan
Forsgren dkk.\ menemukan bahwa tim dengan kinerja pengiriman terbaik
menggabungkan cabangnya ke trunk dalam waktu kurang dari sehari
\cite{forsgren2018accelerate}. Alasannya terlihat ketika dua cabang yang
sama-sama mengubah satu file digabung.

Analoginya seperti dua mahasiswa yang menyunting satu laporan bersama,
masing-masing di \textit{copy} sendiri. Bila disatukan tiap sore, paling
banyak satu dua paragraf yang bertabrakan. Bila baru disatukan dua minggu
kemudian, separuh hari terakhir habis untuk memilih kalimat mana yang menang.

\berkas{simulasi-cabang.py} mengukurnya dengan Git sungguhan. Dua pengembang
masing-masing mengubah tiga baris acak per hari pada katalog 300 baris, lalu
cabang kedua digabung ke cabang pertama, dan blok konfliknya dihitung, seperti
pada Listing \ref{lst:gabung01}.

\begin{lstlisting}[style=PythonStyle, caption={Menggabungkan dua cabang lalu menghitung blok konflik, tersedia di \berkas{sources/chapter01/simulasi-cabang.py}}, label={lst:gabung01}]
  git(direktori, "checkout", "-q", f"a-{umur_hari}")
  hasil = git(direktori, "merge", "--no-edit", f"b-{umur_hari}",
              boleh_gagal=True)
  isi = (direktori / NAMA_FILE).read_text()
  blok_konflik = sum(1 for baris in isi.splitlines()
                     if baris.startswith(PENANDA_KONFLIK))
  if hasil.returncode != 0:
    git(direktori, "merge", "--abort")
  git(direktori, "checkout", "-q", "main")
\end{lstlisting}

Tabel \ref{tab:cabang01} memuat hasilnya. Urutan acaknya sama pada tiap umur,
sehingga cabang tiga hari adalah awal dari cabang tujuh hari, dan yang berbeda
hanya kapan cabangnya digabung.

\begin{table}[!htb]
  \centering
  \caption{Umur cabang lawan konflik ketika digabung, hasil pengukuran nyata
           dengan Git 2.43}
  \label{tab:cabang01}
  \begin{tabularx}{\textwidth}{Xllll}
    \hline
    \textbf{Umur cabang} & \textbf{Diubah A} & \textbf{Diubah B} & \textbf{Diubah keduanya} & \textbf{Blok konflik} \\
    \hline
    1 hari & 3 & 3 & 0 & 0 \\
    \hline
    3 hari & 9 & 9 & 0 & 0 \\
    \hline
    7 hari & 20 & 19 & 1 & 3 \\
    \hline
    14 hari & 40 & 39 & 2 & 11 \\
    \hline
  \end{tabularx}
\end{table}

Dua contoh dari tabel itu:
\begin{enumerate}
  \item Cabang yang digabung sesudah satu atau tiga hari tidak menimbulkan
        konflik sama sekali, walaupun keduanya mengubah file yang sama.
  \item Cabang yang dibiarkan hidup dua minggu menimbulkan 11 blok konflik.
        Umurnya dua kali umur cabang tujuh hari, tetapi konfliknya hampir
        empat kali lipat.
\end{enumerate}

Baris terakhir memuat kejanggalan. Hanya dua baris yang diubah kedua
pengembang, tetapi blok konfliknya sebelas. Git tidak hanya menolak dua
perubahan pada baris yang sama, tetapi juga dua perubahan pada baris yang
berdampingan, karena tidak ada baris tak berubah di antaranya yang dapat
dijadikan patokan \cite{chacon2014progit}. Peluang dua perubahan berdampingan
tumbuh jauh lebih cepat daripada peluang keduanya jatuh di baris yang sama.

Cabang yang pendek baru separuh jalan. Separuh lainnya adalah
\textit{code review}: tiap perubahan dibaca minimal satu orang lain sebelum
digabung. Tabel \ref{tab:review01} memuat aturan yang berlaku sepanjang
semester.

\begin{table}[!htb]
  \centering
  \caption{Aturan \textit{code review} sepanjang semester}
  \label{tab:review01}
  \begin{tabularx}{\textwidth}{XX}
    \hline
    \textbf{Aturan} & \textbf{Alasan} \\
    \hline
    satu perubahan, satu tujuan, sekitar 100 baris & perubahan 1.000 baris umumnya terlalu besar untuk dibaca dengan teliti \\
    \hline
    CI hijau sebelum diminta dibaca & pembaca tidak membuang waktu pada kesalahan yang dapat ditemukan mesin \\
    \hline
    minimal satu pembaca selain penulis & penulis paling sulit melihat kesalahannya sendiri \\
    \hline
    komentar tentang kode, bukan tentang orang & yang diperbaiki kodenya, dan penulisnya tetap mau meminta \textit{review} \\
    \hline
  \end{tabularx}
\end{table}

Batas 100 baris pada baris pertama mengikuti pedoman \textit{review} Google,
yang menyebut 100 baris umumnya ukuran yang wajar dan 1.000 baris umumnya
terlalu besar, dengan keputusan akhir tetap di tangan pembacanya
\cite{google-small-cls}.

\section{Kerangka Continuous Integration}

Continuous Integration berarti tiap perubahan digabung ke cabang utama
sesering mungkin, dan tiap penggabungan diperiksa otomatis dengan menjalankan
seluruh pemeriksaan \cite{humble2010cd}. Pemeriksaan itu disusun sebagai
\textit{pipeline}: rangkaian langkah yang berhenti di langkah pertama yang
gagal.

Analoginya seperti pemeriksa naskah di percetakan kampus. Naskah yang salah
ketik di halaman pertama ditolak sebelum masuk mesin cetak, bukan sesudah 500
eksemplar telanjur tercetak. Karena itu pemeriksaan yang paling murah dan
paling sering menemukan kesalahan diletakkan paling depan.

Gambar \ref{fig:ci01} memperlihatkan tiga langkah CI bab ini. Bab 6
menambahkan analisis statis dan pengujian lewat browser, dan Bab 9
menambahkan pembangunan \textit{image} serta rilisnya.

\begin{figure}[!htb]
  \centering
  \begin{tikzpicture}[
    font=\scriptsize,
    kotak/.style={draw, rounded corners, minimum height=0.85cm,
                  minimum width=2.0cm, align=center, font=\scriptsize\bfseries},
    panah/.style={-{Stealth[length=2mm]}, thick}
  ]
    \node[kotak, fill=blue!8] (commit) at (0,0) {Perubahan\\{\tiny\mdseries satu commit}};
    \node[kotak, fill=blue!8] (sintaks) at (2.6,0) {Periksa\\sintaks};
    \node[kotak, fill=orange!12] (unit) at (5.2,0) {Uji unit\\{\tiny\mdseries 10 uji}};
    \node[kotak, fill=blue!8] (api) at (7.8,0) {Uji API\\{\tiny\mdseries 6 uji}};
    \node[kotak, fill=green!10] (hijau) at (10.4,0) {Hijau\\{\tiny\mdseries boleh digabung}};
    \node[kotak, fill=red!22] (merah) at (5.2,-1.5) {Merah\\{\tiny\mdseries berhenti di sini}};

    \draw[panah] (commit) -- (sintaks);
    \draw[panah] (sintaks) -- (unit);
    \draw[panah] (unit) -- (api);
    \draw[panah] (api) -- (hijau);
    \draw[panah] (unit) -- node[font=\tiny, right] {satu uji gagal} (merah);
  \end{tikzpicture}
  \caption{\textit{Pipeline} CI bab ini: langkah termurah lebih dahulu, dan
           berhenti di langkah pertama yang gagal}
  \label{fig:ci01}
\end{figure}

Langkah-langkah itu ditulis dua kali dengan isi yang sama: sekali sebagai
\textit{workflow} GitHub Actions untuk server CI, dan sekali sebagai skrip
yang dijalankan di komputer sendiri. Listing \ref{lst:ciyml01} memperlihatkan
yang pertama. Filter \texttt{paths} membuat \textit{pipeline} ini hanya
berjalan ketika bagian monorepo miliknya berubah \cite{github-actions-syntax}.

\begin{lstlisting}[language=yaml, caption={\textit{Workflow} CI bab ini, tersedia di \berkas{sources/chapter01/ci.yml}}, label={lst:ciyml01}]
on:
  push:
    branches: [main]
    paths: ["sources/chapter01/**"]
  pull_request:
    paths: ["sources/chapter01/**"]

jobs:
  periksa:
    runs-on: ubuntu-latest
    steps:
      - uses: actions/checkout@v4
      - uses: actions/setup-python@v5
        with:
          python-version: "3.12"

      - name: Sintaks, uji unit, dan uji API
        run: python3 sources/chapter01/ci-lokal.py
\end{lstlisting}

Langkah sesungguhnya tinggal di \berkas{ci-lokal.py}, seperti pada Listing
\ref{lst:langkah01}. Karena server CI dan komputer sendiri menjalankan skrip
yang sama, yang hijau di satu tempat juga hijau di tempat lain, dan kegagalan
di server CI dapat diulang di laptop.

\begin{lstlisting}[style=PythonStyle, caption={Tiga langkah CI, dari yang termurah, tersedia di \berkas{sources/chapter01/ci-lokal.py}}, label={lst:langkah01}]
DAFTAR_LANGKAH = [
    ("Periksa sintaks", ["-m", "compileall", "-q", "."]),
    ("Uji unit service", ["-m", "unittest", "test_layanan"]),
    ("Uji API lewat HTTP", ["-m", "unittest", "test_api"]),
]
\end{lstlisting}

Tabel \ref{tab:ci01} memuat dua kali jalan. Jalan pertama memakai kode apa
adanya. Jalan kedua memakai kode yang sengaja dirusak: aturan harga diubah
sehingga buku seharga nol rupiah ikut diterima.

\begin{table}[!htb]
  \centering
  \caption{CI pada kode utuh dan pada kode yang dirusak, hasil pengukuran
           nyata}
  \label{tab:ci01}
  \begin{tabularx}{\textwidth}{Xllll}
    \hline
    \textbf{Langkah} & \textbf{Utuh} & \textbf{Detik} & \textbf{Dirusak} & \textbf{Detik} \\
    \hline
    periksa sintaks & lolos & 0,04 & lolos & 0,04 \\
    \hline
    uji unit \textit{service} & lolos & 0,06 & gagal & 0,07 \\
    \hline
    uji API lewat HTTP & lolos & 0,10 & tidak dijalankan & - \\
    \hline
    \textbf{hasil} & \textbf{hijau} & & \textbf{merah} & \\
    \hline
  \end{tabularx}
\end{table}

Dua contoh dari tabel itu:
\begin{enumerate}
  \item Kode utuh lolos ketiga langkah dalam 0,20 detik. Pemeriksaan
        semurah itu tidak punya alasan untuk dilewati sebelum commit.
  \item Kode yang dirusak berhenti di uji unit dengan satu nama uji yang
        gagal, \texttt{test\_harga\_nol\_ditolak}. Kesalahannya ditemukan dalam
        0,07 detik, sebelum ada orang yang membacanya, dan jauh sebelum ada
        pembeli yang mendapat buku gratis.
\end{enumerate}

Langkah uji API pada jalan kedua tidak dijalankan, padahal salah satu ujinya
juga akan gagal karena memakai harga nol. Itulah maksud urutan dari yang
termurah: kesalahan yang dapat ditemukan uji unit tidak perlu menunggu server
dinyalakan. CI baru bermakna bila hasil merahnya menahan penggabungan, sehingga
cabang \texttt{main} dilindungi agar hanya menerima perubahan yang hijau.

\section{Kerangka API Minimal: Route, Service, dan Repository}

Praktikum bab-bab berikutnya membutuhkan sistem untuk diukur. Sistem itu
dimulai dari API katalog kecil yang dipecah menjadi tiga lapisan. Lapisan
\textit{route} menerjemahkan permintaan HTTP beserta isinya, yang berbentuk
JavaScript Object Notation (JSON), menjadi panggilan. Lapisan \textit{service}
memegang aturan bisnis. Lapisan \textit{repository} menyimpan data, dan
hanya lapisan itulah yang tahu bentuk penyimpanannya \cite{fowler2002poeaa}.

Analoginya seperti loket pembayaran uang kuliah. Petugas loket menerima
formulir dan memberi tanda terima, tetapi tidak menghitung denda. Bagian
keuangan menghitung denda menurut aturan, tetapi tidak menyimpan berkas.
Arsip menyimpan berkas, tetapi tidak pernah memeriksa aturan. Ketika arsipnya
pindah ke ruangan lain, loket dan bagian keuangan tetap bekerja seperti biasa.

Gambar \ref{fig:lapisan01} memperlihatkan ketiga lapisan beserta dua
\textit{repository} yang dapat dipasang di bawahnya.

\begin{figure}[!htb]
  \centering
  \begin{tikzpicture}[
    font=\scriptsize,
    kotak/.style={draw, rounded corners, minimum height=0.85cm,
                  minimum width=2.4cm, align=center, font=\scriptsize\bfseries},
    panah/.style={-{Stealth[length=2mm]}, thick}
  ]
    \node[kotak, fill=blue!8] (http) at (0,0) {Klien HTTP\\{\tiny\mdseries curl, browser}};
    \node[kotak, fill=blue!8] (rute) at (3.6,0) {Route\\{\tiny\mdseries rute.py}};
    \node[kotak, fill=orange!12] (layanan) at (7.4,0) {Service\\{\tiny\mdseries layanan.py}};
    \node[kotak, fill=green!10] (memori) at (5.8,-1.6) {Repository memori\\{\tiny\mdseries uji unit}};
    \node[kotak, fill=green!10] (sqlite) at (9.0,-1.6) {Repository SQLite\\{\tiny\mdseries aplikasi}};

    \draw[panah] (http) -- node[font=\tiny, above] {JSON} (rute);
    \draw[panah] (rute) -- node[font=\tiny, above] {panggilan} (layanan);
    \draw[panah] (layanan.south) -- (memori.north);
    \draw[panah] (layanan.south) -- (sqlite.north);
    \node[font=\tiny, gray] at (7.4,0.75) {aturan bisnis};
  \end{tikzpicture}
  \caption{Tiga lapisan API, dengan \textit{repository} yang dapat diganti}
  \label{fig:lapisan01}
\end{figure}

Aturan bisnisnya tertulis pada Listing \ref{lst:layanan01}. Tidak ada satu
pun baris HTTP atau Structured Query Language (SQL) di sana, sehingga aturan ini dapat diuji tanpa server
dan tanpa basis data.

\begin{lstlisting}[style=PythonStyle, caption={Aturan buku yang sah di lapisan \textit{service}, tersedia di \berkas{sources/chapter01/layanan.py}}, label={lst:layanan01}]
  def tambah_buku(self, judul, harga, stok):
    """Memeriksa satu buku baru terhadap aturan katalog, lalu menyimpannya."""
    judul = judul.strip() if isinstance(judul, str) else ""
    if not judul:
      raise KesalahanValidasi("judul wajib diisi")
    if len(judul) > PANJANG_JUDUL_MAKSIMUM:
      raise KesalahanValidasi(
          f"judul paling panjang {PANJANG_JUDUL_MAKSIMUM} karakter")
    if not bilangan_bulat(harga) or harga < HARGA_MINIMUM_RUPIAH:
      raise KesalahanValidasi("harga harus bilangan bulat positif")
    if not bilangan_bulat(stok) or stok < STOK_MINIMUM:
      raise KesalahanValidasi("stok harus bilangan bulat, tidak negatif")
    return self.repositori.simpan(judul, harga, stok)
\end{lstlisting}

Kesalahan dari \textit{service} diterjemahkan \textit{route} menjadi status
HTTP. Listing \ref{lst:rute01} memperlihatkan satu terjemahan itu:
\texttt{BukuTidakDitemukan} menjadi 404.

\begin{lstlisting}[style=PythonStyle, caption={\textit{Route} menerjemahkan kesalahan menjadi status HTTP, tersedia di \berkas{sources/chapter01/rute.py}}, label={lst:rute01}]
    nomor = self.path.removeprefix(AWALAN_BUKU)
    if self.path.startswith(AWALAN_BUKU) and nomor.isdigit():
      try:
        self.kirim_json(200, layanan.cari(int(nomor)))
      except BukuTidakDitemukan as kesalahan:
        self.kirim_json(404, {"pesan": str(kesalahan)})
      return
    self.kirim_json(404, {"pesan": "alamat tidak dikenal"})
\end{lstlisting}

Pemisahan itu terlihat paling jelas pada pengujiannya. Tabel
\ref{tab:uji01} membandingkan uji di lapisan \textit{service} dengan uji
lewat HTTP pada aplikasi yang sama.

\begin{table}[!htb]
  \centering
  \caption{Uji di lapisan \textit{service} lawan uji lewat HTTP, hasil
           pengukuran nyata}
  \label{tab:uji01}
  \begin{tabularx}{\textwidth}{Xlll}
    \hline
    \textbf{Lapisan yang diuji} & \textbf{Jumlah uji} & \textbf{Detik} & \textbf{Yang dibutuhkan} \\
    \hline
    \textit{service}, \textit{repository} memori & 10 & di bawah 0,001 & tidak ada \\
    \hline
    API lewat HTTP, SQLite & 6 & 0,016 & server dan basis data \\
    \hline
  \end{tabularx}
\end{table}

Tabel itu punya dua konsekuensi praktis. Pertama, aturan bisnis diuji di
lapisan \textit{service}, karena di sanalah uji paling murah dan paling banyak
dapat ditulis. Uji lewat HTTP cukup memeriksa hal yang memang hanya ada di
HTTP: status, bentuk JSON, dan terjemahan
kesalahan. Kedua, \textit{repository} dapat diganti tanpa menyentuh dua
lapisan di atasnya. Bab 3 mengganti SQLite dengan PostgreSQL, dan
\textit{service} serta \textit{route}-nya tidak perlu berubah.

Angka API juga sempat menipu. Pada percobaan pertama, keenam uji API memakan
0,506 detik, hampir seluruhnya menunggu server berhenti, karena server
pustaka standar Python baru memeriksa perintah berhenti tiap setengah detik.
Selang itu diperkecil di \berkas{test\_api.py}, dan angkanya turun menjadi
0,016 detik. Uji yang lambat tanpa alasan akhirnya dilewati orang, sehingga
lamanya pun perlu diukur.

\section{Ringkasan}

Mata kuliah ini menilai aplikasi web dari sifat yang tidak terlihat di layar:
cepatnya, ketahanannya, keamanannya, dan kemudahannya diubah. Jawabannya
selalu angka yang diukur sendiri, misalnya 199,8 lawan 1,8 milidetik untuk
peristiwa yang sama pada Bab 13. Keempat belas pertemuannya dibagi menjadi
empat bagian, dan urutannya disusun untuk menopang enam topik penelitian.

Lingkungan kerja diperiksa sebelum praktikum pertama dengan satu
\textit{smoke test}. Pada laptop penulisan modul ini, seluruh perkakas lolos,
tetapi tiga dari 37 port sudah dipakai proyek lain, yaitu port milik Bab 3,
Bab 4, dan Bab 11. Kegagalan seperti itu jauh lebih murah ditemukan pada
pertemuan pertama daripada di tengah praktikum bab yang membutuhkannya.

Alur kerjanya memakai satu monorepo, cabang berumur pendek, dan
\textit{code review}. Pengukuran dengan Git memperlihatkan alasannya: cabang
yang digabung sesudah satu atau tiga hari tidak menimbulkan konflik,
sedangkan cabang dua minggu menimbulkan 11 blok konflik walaupun hanya dua
baris yang diubah kedua pihak, karena perubahan pada baris yang berdampingan
ikut bertabrakan.

CI menjalankan pemeriksaan termurah lebih dahulu dan berhenti di langkah
pertama yang gagal. Kode utuh lolos dalam 0,20 detik, sedangkan kode yang
dirusak tertahan di uji unit dalam 0,07 detik dengan nama uji yang gagal.
Kerangka API berpola \textit{route}, \textit{service}, dan \textit{repository}
membuat aturan bisnis dapat diuji tanpa server: sepuluh uji selesai di bawah
satu milidetik.

Pesan utama bab ini adalah bahwa kebiasaan kerja ditetapkan sejak awal dan
dijaga mesin. Lingkungan diperiksa \textit{smoke test}, cabang digabung
selagi kecil, dan perubahan cacat ditolak CI sebelum sempat dibaca orang,
sehingga tiga belas bab berikutnya dapat dipakai untuk mengukur, bukan untuk
membereskan.

\ifajar
  \ifajar
  \section{Latihan}
\else
  \section{Bab 1: Overview Mata Kuliah dan Alur Rekayasa}
\fi

Seluruh file yang dipakai pada latihan ini tersedia di
\berkas{sources/chapter01/}, dihitung relatif terhadap akar repositori.
Latihan ini hanya membutuhkan Python 3 dan Git. Docker, Node.js, dan Chrome
hanya diperiksa keberadaannya.

Seluruh skrip hanya menyentuh komputer sendiri. Pemeriksaan port mencoba
memesan tiap port sebentar lalu melepasnya, tanpa mengirim data ke mana pun,
dan percobaan Git berjalan di direktori sementara. Angka detik tiap mahasiswa
akan berbeda, sedangkan pola perbandingannya akan sama.

\ifajar
  \subsection*{Latihan 1: Smoke Test Lingkungan}
\else
  \subsection*{Prosedur 1: Smoke Test Lingkungan}
\fi

Latihan ini memastikan seluruh perkakas keempat belas bab sudah siap, dan
memperlihatkan bahwa kegagalan yang ditemukan sekarang jauh lebih murah
daripada yang ditemukan di tengah praktikum.

\begin{enumerate}
  \item Buat venv bersama bila belum ada, lalu pasang pustaka basis data,
        seperti pada Listing \ref{lst:c01-lat1-jalan}.
  \item Jalankan \berkas{periksa-lingkungan.py}, lalu catat seluruh barisnya.
  \item Perbaiki tiap pemeriksaan yang gagal: pasang perkakas yang kurang,
        nyalakan Docker, atau hentikan proses yang memakai port bab lain.
  \item Jalankan ulang sampai hasilnya hijau, lalu catat apa saja yang
        diperbaiki.
\ifajar
  \item Isi \tabisi{tab:c01-lat1-hasil}. Tabel \ref{tab:c01-lat1-contoh}
        memperlihatkan contoh yang sudah terisi.
\else
  \item Bandingkan hasil pengukuran sendiri dengan Tabel
        \ref{tab:c01-lat1-contoh}.
\fi
  \item Jawab tiga pertanyaan berikut. Pertama, mengapa pemeriksaan Docker
        menanyai server, bukan hanya mencari perintah \texttt{docker}?
        Kedua, bab mana yang tertahan bila port 8080 sudah terpakai, dan
        kapan kegagalannya akan muncul tanpa \textit{smoke test} ini?
        Ketiga, mengapa versi minimum ditulis di skrip, bukan hanya di
        naskah?
\end{enumerate}

\begin{lstlisting}[language=bash, caption={Menyiapkan venv lalu menjalankan \textit{smoke test}, tersedia di \berkas{sources/chapter01/periksa-lingkungan.py}}, label={lst:c01-lat1-jalan}]
cd sources/chapter01
python3 -m venv ../.venv
../.venv/bin/pip install "psycopg[binary,pool]"
# Dua perintah di atas cukup sekali, bila sources/.venv belum ada.
python3 periksa-lingkungan.py

# Keluarannya periksa-lingkungan.py:
#   Pemeriksaan                      Hasil  Keterangan
#   Python 3.11 atau lebih baru      lolos  3.15.2
#   Node.js 20 atau lebih baru       lolos  v22.23.4
#   npm 10 atau lebih baru           lolos  10.9.9
#   Git 2.30 atau lebih baru         lolos  git version 2.43.0
#   Docker 24 atau lebih baru        lolos  server 29.9.3
#   Docker Compose 2 atau lebih baru lolos  5.7.0
#   Chrome atau Chromium             lolos  Google Chrome 157.0.8103.38
#   psycopg di sources/.venv         lolos  3.3.5
#   Port bab 1 sampai 14             gagal  3 terpakai: 5433, 8000, 8080
#     port 5433 milik Bab 3 sedang dipakai proses lain
#     port 8000 milik Bab 4 sedang dipakai proses lain
#     port 8080 milik Bab 11 sedang dipakai proses lain
#
#   Hasil: merah, 1 pemeriksaan gagal
\end{lstlisting}

\begin{table}[!htb]
  \centering
  \caption{Contoh hasil Latihan 1 yang sudah terisi}
  \label{tab:c01-lat1-contoh}
  \begin{tabularx}{\textwidth}{Xll}
    \hline
    \textbf{Pemeriksaan} & \textbf{Hasil} & \textbf{Keterangan} \\
    \hline
    perkakas dan versinya, tujuh baris & lolos & seluruhnya di atas minimum \\
    \hline
    psycopg di venv & lolos & 3.3.5 \\
    \hline
    port bab 1 sampai 14 & gagal & 5433, 8000, 8080 \\
    \hline
    \textbf{hasil} & \textbf{merah} & 1 pemeriksaan gagal \\
    \hline
  \end{tabularx}
\end{table}

\ifajar
\begin{table}[!htb]
  \centering
  \caption{Lembar isian hasil pemeriksaan Latihan 1}
  \label{tab:c01-lat1-hasil}
  \begin{tabularx}{\textwidth}{Xll}
    \hline
    \textbf{Pemeriksaan} & \textbf{Hasil} & \textbf{Keterangan} \\
    \hline
    perkakas dan versinya, tujuh baris & \dots & \dots \\
    \hline
    psycopg di venv & \dots & \dots \\
    \hline
    port bab 1 sampai 14 & \dots & \dots \\
    \hline
    \textbf{hasil} & \dots & \dots \\
    \hline
  \end{tabularx}
\end{table}
\fi

Contoh di atas merah karena ketiga port dipakai proyek lain di laptop yang
sama, bukan karena perkakasnya kurang. Hasil di laptop lain hampir pasti
berbeda. Yang dilaporkan adalah pemeriksaan mana yang gagal dan apa yang
dilakukan sampai hasilnya hijau.

\ifajar
  \subsection*{Latihan 2: Umur Cabang dan Konflik}
\else
  \subsection*{Prosedur 2: Umur Cabang dan Konflik}
\fi

Latihan ini membuktikan bahwa konflik tumbuh lebih cepat daripada umur
cabangnya, sehingga cabang pantas digabung selagi pendek.

\begin{enumerate}
  \item Jalankan \berkas{simulasi-cabang.py}, seperti pada Listing
        \ref{lst:c01-lat2-jalan}, lalu catat keempat barisnya.
  \item Ubah \texttt{BARIS\_PER\_HARI} di \berkas{simulasi-cabang.py} menjadi
        6, jalankan ulang, lalu catat kolom blok konfliknya.
  \item Kembalikan \texttt{BARIS\_PER\_HARI} menjadi 3.
\ifajar
  \item Isi \tabisi{tab:c01-lat2-hasil}. Tabel \ref{tab:c01-lat2-contoh}
        memperlihatkan contoh yang sudah terisi.
\else
  \item Bandingkan hasil pengukuran sendiri dengan Tabel
        \ref{tab:c01-lat2-contoh}.
\fi
  \item Jawab tiga pertanyaan berikut. Pertama, mengapa blok konflik pada 14
        hari jauh lebih banyak daripada baris yang diubah kedua pengembang?
        Kedua, apa yang terjadi pada konflik ketika tiap pengembang bekerja
        dua kali lebih banyak per hari? Ketiga, berapa hari umur cabang
        terpanjang yang masih bebas konflik pada kedua setelan?
\end{enumerate}

\begin{lstlisting}[language=bash, caption={Mengukur konflik menurut umur cabang, tersedia di \berkas{sources/chapter01/simulasi-cabang.py}}, label={lst:c01-lat2-jalan}]
python3 simulasi-cabang.py

# Keluarannya:
#   300 baris, tiap pengembang mengubah 3 baris per hari, git version 2.43.0
#
#   Umur cabang  Diubah A  Diubah B  Diubah keduanya  Blok konflik
#        1 hari         3         3                0             0
#        3 hari         9         9                0             0
#        7 hari        20        19                1             3
#       14 hari        40        39                2            11
#
# Kolom blok konflik dengan BARIS_PER_HARI = 6: 0, 2, 10, 21
\end{lstlisting}

\begin{table}[!htb]
  \centering
  \caption{Contoh hasil Latihan 2 yang sudah terisi}
  \label{tab:c01-lat2-contoh}
  \begin{tabularx}{\textwidth}{Xll}
    \hline
    \textbf{Umur cabang} & \textbf{Konflik, 3 baris per hari} & \textbf{Konflik, 6 baris per hari} \\
    \hline
    1 hari & 0 & 0 \\
    \hline
    3 hari & 0 & 2 \\
    \hline
    7 hari & 3 & 10 \\
    \hline
    14 hari & 11 & 21 \\
    \hline
  \end{tabularx}
\end{table}

\ifajar
\begin{table}[!htb]
  \centering
  \caption{Lembar isian hasil pengukuran Latihan 2}
  \label{tab:c01-lat2-hasil}
  \begin{tabularx}{\textwidth}{Xll}
    \hline
    \textbf{Umur cabang} & \textbf{Konflik, 3 baris per hari} & \textbf{Konflik, 6 baris per hari} \\
    \hline
    1 hari & \dots & \dots \\
    \hline
    3 hari & \dots & \dots \\
    \hline
    7 hari & \dots & \dots \\
    \hline
    14 hari & \dots & \dots \\
    \hline
  \end{tabularx}
\end{table}
\fi

Benih acaknya tetap, sehingga angka latihan ini sama persis pada komputer
mana pun dengan Git versi yang sama. Versi Git yang berbeda dapat menghitung
batas blok konflik sedikit berbeda, sedangkan polanya tetap: konflik tumbuh
jauh lebih cepat daripada umur cabangnya.

\ifajar
  \subsection*{Latihan 3: Kerangka API dan CI yang Hijau}
\else
  \subsection*{Prosedur 3: Kerangka API dan CI yang Hijau}
\fi

Latihan ini menyalakan kerangka API yang dipakai bab-bab berikutnya, lalu
membuktikan bahwa CI menahan perubahan yang melanggar aturan bisnis.

\begin{enumerate}
  \item Jalankan \berkas{ci-lokal.py}, lalu catat ketiga langkahnya beserta
        detiknya.
  \item Nyalakan \berkas{aplikasi.py}, lalu kirim tiga permintaan seperti
        pada Listing \ref{lst:c01-lat3-jalan}. Catat status HTTP tiap
        permintaan.
  \item Ubah \texttt{harga < HARGA\_MINIMUM\_RUPIAH} di \berkas{layanan.py}
        menjadi \texttt{harga < 0}, jalankan \berkas{ci-lokal.py}, lalu catat
        langkah tempatnya berhenti dan nama uji yang gagal.
  \item Kembalikan \berkas{layanan.py}, lalu pastikan CI kembali hijau.
  \item Buat repositori tim, salin direktori ini ke dalamnya, salin
        \berkas{ci.yml} ke \berkas{.github/workflows/}, sesuaikan
        \texttt{paths}-nya, lalu commit.
\ifajar
  \item Isi \tabisi{tab:c01-lat3-hasil}. Tabel \ref{tab:c01-lat3-contoh}
        memperlihatkan contoh yang sudah terisi.
\else
  \item Bandingkan hasil pengukuran sendiri dengan Tabel
        \ref{tab:c01-lat3-contoh}.
\fi
  \item Jawab tiga pertanyaan berikut. Pertama, mengapa uji API tidak
        dijalankan pada jalan yang merah, padahal salah satu ujinya juga
        akan gagal? Kedua, lapisan mana yang berubah bila SQLite diganti
        PostgreSQL? Ketiga, mengapa satu skrip CI yang sama dipakai di
        laptop dan di server CI?
\end{enumerate}

\begin{lstlisting}[language=bash, caption={CI lokal lalu tiga permintaan ke API, tersedia di \berkas{sources/chapter01/}}, label={lst:c01-lat3-jalan}]
python3 ci-lokal.py

# Keluarannya:
#   Langkah              Hasil   Detik  Keterangan
#   Periksa sintaks      lolos    0.04  seluruh file
#   Uji unit service     lolos    0.06  10 uji
#   Uji API lewat HTTP   lolos    0.10  6 uji
#
#   Hasil: hijau, seluruh langkah lolos

python3 aplikasi.py &
curl -s -X POST http://127.0.0.1:8050/buku \
  -H 'Content-Type: application/json' \
  -d '{"judul": "Basis Data Terapan", "harga": 85000, "stok": 4}'
curl -s -X POST http://127.0.0.1:8050/buku \
  -H 'Content-Type: application/json' \
  -d '{"judul": "Struktur Data", "harga": 0, "stok": 2}'
curl -s http://127.0.0.1:8050/buku/9

# Keluarannya, dengan status HTTP di ujung tiap baris:
# {"nomor": 1, "judul": "Basis Data Terapan", "harga": 85000, "stok": 4} 201
# {"pesan": "harga harus bilangan bulat positif"} 400
# {"pesan": "buku nomor 9 tidak ada"} 404
\end{lstlisting}

\begin{table}[!htb]
  \centering
  \caption{Contoh hasil Latihan 3 yang sudah terisi}
  \label{tab:c01-lat3-contoh}
  \begin{tabularx}{\textwidth}{Xll}
    \hline
    \textbf{Pengamatan} & \textbf{Kode utuh} & \textbf{Kode dirusak} \\
    \hline
    hasil CI & hijau, 0,20 detik & merah, berhenti di uji unit \\
    \hline
    uji yang gagal & tidak ada & \texttt{test\_harga\_nol\_ditolak} \\
    \hline
    status buku sah, harga nol, nomor 9 & 201, 400, 404 & 201, 201, 404 \\
    \hline
  \end{tabularx}
\end{table}

\ifajar
\begin{table}[!htb]
  \centering
  \caption{Lembar isian hasil pengamatan Latihan 3}
  \label{tab:c01-lat3-hasil}
  \begin{tabularx}{\textwidth}{Xll}
    \hline
    \textbf{Pengamatan} & \textbf{Kode utuh} & \textbf{Kode dirusak} \\
    \hline
    hasil CI & \dots & \dots \\
    \hline
    uji yang gagal & \dots & \dots \\
    \hline
    status buku sah, harga nol, nomor 9 & \dots & \dots \\
    \hline
  \end{tabularx}
\end{table}
\fi

Kolom kode dirusak pada baris terakhir menunjukkan akibat yang ditahan CI.
Pada kode yang dirusak, permintaan buku seharga nol dijawab 201, artinya buku
gratis benar-benar masuk katalog. Karena itu \textit{route} dan
\textit{repository} tidak cukup diuji sendiri-sendiri: aturan yang salah di
\textit{service} tetap menghasilkan HTTP yang tampak sehat.

\fi
